\documentclass[10pt]{article}
\usepackage{fontspec}
\usepackage[english]{babel}
\usepackage{amsmath,amssymb,mathtools}
\usepackage{geometry}
\usepackage{enumitem}
\usepackage{tabularx}
\usepackage{array,booktabs}
\geometry{margin=0.86in}
\setlength{\parindent}{1.25em}
\setlength{\parskip}{0.14em}
\emergencystretch=3em
\frenchspacing
\setlist[enumerate]{leftmargin=2.2em,itemsep=0.12em,topsep=0.25em}
\newcommand{\A}{\mathbb A}
\newcommand{\C}{\mathbb C}
\newcommand{\D}{\mathbb D}
\newcommand{\N}{\mathbb N}
\newcommand{\Q}{\mathbb Q}
\newcommand{\R}{\mathbb R}
\newcommand{\Z}{\mathbb Z}
\newcommand{\Ql}{\mathbb Q_\ell}
\newcommand{\Zl}{\mathbb Z_\ell}
\newcommand{\Gm}{\mathbb G_m}
\newcommand{\Gr}{\operatorname{Gr}}
\newcommand{\Gal}{\operatorname{Gal}}
\newcommand{\Spec}{\operatorname{Spec}}
\newcommand{\Hom}{\operatorname{Hom}}
\newcommand{\Ker}{\operatorname{Ker}}
\newcommand{\Id}{\operatorname{Id}}
\newcommand{\pr}{\operatorname{pr}}
\newcommand{\card}{\#}
\newcommand{\barK}{\overline{K}}
\newcommand{\wt}{\widetilde}
\newcommand{\calH}{\mathcal H}
\begin{document}
% ---- all ----
\begin{center}
{\small Proceedings, Intern. Math. Congress, 1970. Volume 1, p. 425 to 430.}\\[4em]
{\Large\bfseries HODGE THEORY I}\\[1.6em]
{\large by PIERRE DELIGNE}
\end{center}

We propose to give a heuristic dictionary between statements in $l$-adic cohomology and statements in Hodge theory. This dictionary draws in particular on [3] and on Grothendieck's conjectural theory of motives [2]. So far, it has mainly served to formulate conjectures in Hodge theory, and it has sometimes suggested a proof of them.

\paragraph{Definition 1.1.} — A mixed Hodge structure $H$ consists of
\begin{enumerate}[label=(\alph*)]
\item A finitely generated $\Z$-module $H_\Z$ (the “integral lattice”);
\item A finite growth filtration $W$ on $H_\Q=H_\Z\otimes\Q$ (the “weight filtration”); % [sic] printed croissance; intended croissante.
\item A finite decreasing filtration $F$ on $H_\C=H_\Z\otimes\C$ (the “Hodge filtration”).
\end{enumerate}
These data are subject to the axiom:

There exists on $\Gr_W(H_\C)$ a (unique) bigrading by subspaces $H^{p,q}$ such that
\begin{enumerate}[label=(\roman*)]
\item
\[
\Gr_W^n(H_\C)=\bigoplus_{p+q=n}H^{p,q}
\]
\item the filtration $F$ induces on $\Gr_W(H_\C)$ the filtration
\[
\Gr_W(F)^p=\bigoplus_{p'\ge p}H^{p',q'}
\]
\item $\overline{H^{pq}}=H^{qp}$.
\end{enumerate}
A morphism $f:H\to H'$ is a homomorphism $f_\Z:H_\Z\to H'_\Z$ such that $f_\Q:H_\Q\to H'_\Q$ and $f_\C:H_\C\to H'_\C$ are respectively compatible with the filtrations $W$ and $F$.

The Hodge numbers of $H$ are the integers
\[
\tag{1.2}
h^{pq}=\dim H^{pq}=h^{qp}.
\]
One says that $H$ is pure of weight $n$ if $h^{pq}=0$ for $p+q\ne n$ (i. e. if $\Gr_W^i(H)=0$ for $i\ne n$). One also says that $H$ is a Hodge structure of weight $n$.

The Tate Hodge structure $\Z(1)$ is the Hodge structure of weight $-2$, purely of type $(-1,-1)$, for which $\Z(1)_\C=\C$ and $\Z(1)_\Z=2\pi i\Z=\Ker(\exp:\C\to\C^*)\subset\C$. We set $\Z(n)=\Z(1)^{\otimes n}$.

One can show that mixed Hodge structures form an abelian category. If $f:H\to H'$ is a morphism, then $f_\Q$ and $f_\C$ are strictly compatible with the filtrations $W$ and $F$ ([1], 2.3.5).

\paragraph{2.} Let $A$ be a normal integral domain of finite type over $\Z$, $K$ its field of fractions and $\barK$ an algebraic closure of $K$. Let $K_{nr}$ be the largest subextension of $\barK$ unramified at every prime ideal of $A$. It is known, or taken as a definition, that
\[
\pi_1(\Spec(A),\barK)=\Gal(K_{nr}/K).
\]
For each closed point $x$ of $\Spec(A)$, defined by a maximal ideal $m_x$ of $A$, the residue field $k_x=A/m_x$ is finite; the point $x$ defines a conjugacy class of “Frobenius substitutions” $\varphi_x\in\pi_1(\Spec(A),\barK)$. We set $q_x=\#k_x$ and $F_x=\varphi_x^{-1}$.

Let $K$ be a field of finite type over the prime field of characteristic $p$, $\barK$ an algebraic closure of $K$, $l$ a prime number $\ne p$ and $H$ a finitely generated $\Z_l$- (or $\Q_l$-) module equipped with a continuous action $\rho$ of $\Gal(\barK/K)$. In what follows, we shall always assume that there exists $A$ as above, with $l$ invertible in $A$, such that $\rho$ factors through
\[
\pi_1(\Spec(A),\barK)=\Gal(K_{nr}/K).
\]
We shall say that $H$ is pure of weight $n$ if for every closed point $x$ of a nonempty open subset of $\Spec(A)$, the eigenvalues $\alpha$ of $F_x$ acting on $H$ are algebraic integers all of whose complex conjugates have absolute value
\[
|\alpha|=q_x^{n/2}.
\]

\paragraph{Principle 2.1.} — If the Galois module $H$ “comes from algebraic geometry”, there exists on $H_{\Q_l}=H\otimes_{\Z_l}\Q_l$ a (unique) increasing filtration $W$ (the “weight filtration”), invariant under Galois, such that $\Gr_n^W(H)$ is pure of weight $n$.

One may expect that $\Gr_n^W(H)$ is moreover semisimple.

When resolution of singularities is available, one can often give a conjectural definition of $W$, whose correctness follows from the Weil conjectures [5] (cf. 6).

Let $\mu$ be the subgroup of $\barK^*$ consisting of the roots of unity. The Tate module $\Z_l(1)$, defined by
\[
\Z_l(1)=\Hom(\Q_l/\Z_l,\mu)
\]
is pure of weight $-2$. We set $\Z_l(n)=\Z_l(1)^{\otimes n}$.

It is trivial that every morphism $f:H\to H'$ is strictly compatible with the weight filtration.

Principle 2.1 agrees with the fact that every extension of $\Gm$ (“weight $-2$”) by an abelian variety (“weight $-1>-2$”) is trivial.

\paragraph{3. Translation.} — The Galois modules that appear in $l$-adic cohomology have mixed Hodge structures as their analogues over $\C$. We also have the dictionary
\[
\begin{array}{r@{\quad\vrule\quad}l}
\text{pure module of weight }n & \text{Hodge structure of weight }n\\
\text{weight filtration} & \text{weight filtration}\\
\text{Galois-equivariant homomorphism} & \text{morphism}\\
\text{Tate module }\Z_l(1) & \text{Tate Hodge structure }\Z(1)
\end{array}
\]

\paragraph{4.} Let $X$ be a complex algebraic variety (= a scheme of finite type over $\C$, which we shall assume to be separated). There exists a subfield $K$ of $\C$, of finite type over $\Q$, such that $X$ can be defined over $K$ (i. e. arises by extension of scalars from $K$ to $\C$ from a $K$-scheme $X'$). Let $\barK$ be the algebraic closure of $K$ in $\C$. The Galois group $\Gal(\barK/K)$ then acts on the $l$-adic cohomology groups $H^*(X,\Z_l)$; we have
\[
H^*(X(\C),\Z)\otimes\Z_l=H^*(X,\Z_l)=H^*(X'_{\barK},\Z_l).
\]

According to 3, there is reason to expect the cohomology groups $H^n(X(\C),\Z)$ to carry natural mixed Hodge structures. This can be proved (see [1], 3.2.5, for the case where $X$ is smooth; the proof is algebraic, starting from classical Hodge theory [6]). For $X$ projective and smooth, the Weil conjectures imply that $H^n(X,\Z_l)$ is pure of weight $n$, while classical Hodge theory equips $H^n(X,\Z)$ with a Hodge structure of weight $n$. For every morphism $f:X\to Y$ and for $K$ large enough, $f^*:H^*(Y,\Z_l)\to H^*(X,\Z_l)$ commutes with Galois (by transport of structure); similarly, $f^*:H^*(Y,\Z)\to H^*(X,\Z)$ is a morphism of mixed Hodge structures. For $X$ smooth, the cohomology class in $H^{2n}(X,\Z_l(n))$ of an algebraic cycle of codimension $n$, $Z$, defined over $K$, is invariant under Galois, i. e. defines
\[
c(Z)\in\Hom_{\Gal}(\Z_l(-n),H^{2n}(X,\Z_l)).
\]
Similarly, the cohomology class $c(Z)\in H^{2n}(X(\C),\Z)$ is purely of type $(n,n)$, i. e. corresponds to
\[
c(Z)\in\Hom_{H.M.}(\Z(-n),H^{2n}(X(\C),\Z)).
\]

\paragraph{5.} If $f:H\to H'$ is a Galois-equivariant morphism between $\Q_l$-vector spaces of different weights, then $f=0$. Similarly, if $f:H\to H'$ is a morphism of pure mixed Hodge structures of different weights, then $f$ is torsion. A more useful remark is the % [sic] printed "Une remarque plus utile est la"; intended "Une remarque plus utile est la suivante.".

\paragraph{Scholium 5.1.} — Let $H$ and $H'$ be Hodge structures of weights $n$ and $n'$, with $n>n'$. Let $f:H_\Q\to H'_\Q$ be a homomorphism such that $f:H_\C\to H'_\C$ respects $F$. Then $f=0$.

\paragraph{6.} Let $X$ be a smooth projective variety over $\C$, $D=\sum_1^n D_i$ a divisor with normal crossings in $X$, a sum of smooth divisors, and $j$ the inclusion into $X$ of $U=X-D$. For $Q\subset[1,n]$, we set
\[
D_Q=\bigcap_{i\in Q}D_i.
\]
In $l$-adic cohomology, we have canonically
\[
\tag{6.1}
R^qj_*\Z_l=\bigoplus_{\#Q=q}\Z_l(-q)_{D_Q},
\]
and the Leray spectral sequence for $j$ takes the form
\[
\tag{6.2}
E_2^{pq}=\bigoplus_{\#Q=q}H^p(D_Q,\Q_l)\otimes\Z_l(-q)\Rightarrow H^{p+q}(U,\Q_l).
\]
According to the Weil conjectures [5], $H^p(D_Q,\Q_l)$ is pure of weight $p$, so that $E_2^{pq}$ is pure of weight $p+2q$. As a quotient of a subobject of $E_2^{pq}$, $E_r^{pq}$ is also pure of weight $p+2q$. By 5, $d_r=0$ for $r\ge3$, since the weights $p+2q$ and $p+2q-r+2$ of $E_r^{pq}$ and $E_r^{p+r,q-r+1}$ are different. We therefore have $E_3^{pq}=E_\infty^{pq}$. Up to reindexing, the weight filtration of $H^*(U,\Q_l)$ is the abutment of (6.2)
\[
\tag{6.3}
\Gr_n^W(H^k(U,\Q_l))=E_3^{2k-n,n-k}
\]

\paragraph{7.} In integral cohomology, for the usual topology, the Leray spectral sequence for $j$ takes the form
\[
\tag{7.1}
{}'E_2^{pq}=\bigoplus_{\#Q=q}H^p(D_Q,\Z)\Rightarrow H^{p+q}(U,\Z).
\]
Since each $D_Q$ is a nonsingular projective variety, ${}'E_2^{pq}$ is equipped with a Hodge structure of weight $p$. We set $E_2^{pq}={}'E_2^{pq}\otimes\Z(-q)$ (a Hodge structure of weight $p+2q$). As an abelian group, $E_2^{pq}={}'E_2^{pq}$; it is useful to regard (7.1) as a spectral sequence with initial term $E_2^{pq}$. According to 3, one should expect
\[
d_2:E_2^{pq}\to E_2^{p+2,q-1}
\]
to be a morphism of Hodge structures. This is proved by interpreting $d_2$ as a Gysin map. It follows that $E_3^{pq}$ is equipped with a Hodge structure of weight $p+2q$. According to 3, one expects that, modulo torsion, the spectral sequence (6.4) % [sic] printed (6.4); intended (7.1).
degenerates at the $E_3$ term ($E_3=E_\infty$), and that the vanishing of the $d_r$ ($r\ge3$) is an application of 5.1. This program is carried out in [1] 3.2. There the weight filtration of $H^*(U,\Q)$ is defined as the abutment of (7.1), up to the reindexing (6.3).

In fact, to equip cohomology groups $H^*$ with a mixed Hodge structure, the key point so far has always been to find a spectral sequence $E$ with abutment $H^*$ such that the $l$-adic analogue of $E_2^{pq}$ is conjecturally pure (of weight $p+2q$); $E_2^{pq}$ must then carry a natural Hodge structure (of weight $p+2q$), and the filtration $W$ is the abutment of $E$.

\paragraph{8.} Let $\Spec(V)$ be the spectrum of a henselian discrete valuation ring (a henselian trait) with field of fractions $K$ and residue field $k$ of finite type over the prime field of characteristic $p$. Let $\barK$ be an algebraic closure of $K$ and $H$ a finite-dimensional vector space over $\Q_l$ ($l\ne p$), on which $\Gal(\barK/K)$ acts continuously. Grothendieck has shown ([4], appendix) that a subgroup of finite index in the inertia group $I$ acts unipotently. Replacing $V$ by a finite extension, we reduce to the case where the action of all of $I$ is unipotent (the semistable case); it then factors the largest pro-$l$-group quotient $I_l$ of $I$, canonically isomorphic to $\Z_l(1)$. % [sic] printed "elle se factorise alors le plus grand"; intended "elle se factorise alors par le plus grand".

\paragraph{Principle 8.1.} — In the semistable case, if the Galois module $H$ “comes from algebraic geometry”, there exists a (unique) increasing filtration $W$ of $H$ (the “weight filtration”) such that $I$ acts trivially on $\Gr_n^W(H)$ and such that $\Gr_n^W(H)$, as a Galois module under $\Gal(\bar k/k)\simeq\Gal(\barK/K)/I$, is pure of weight $n$.

Compare with 2.1 and with the appendix of [4].

When resolution of singularities is available, one can sometimes give a conjectural definition of $W$, whose validity follows from the Weil conjectures. With the help of resolution and Weil, it is often easy to show that in any case $H$ has a filtration with pure Galois modules (under $\Gal(\bar k/k)$) as successive quotients.

Suppose $H$ is semistable. For $T\in I_l$, we define $\log T$ as the finite sum
\[
-\sum_{n>0}(\mathit{Id}-T)^n/n.
\]
The map $(T,x)\to\log T(x)$ is identified with a homomorphism
\[
\tag{8.2}
M:\Z_l(1)\otimes H\longrightarrow H.
\]
Since $\Z_l(1)$ has weight $-2$, we necessarily have (cf. 5)
\[
\tag{8.3}
M(\Z_l(1)\otimes W_n(H))\subset W_{n-2}(H)
\]
and $M$ induces
\[
\tag{8.4}
\Gr(M):\Z_l(1)\otimes\Gr_n^W(H)\longrightarrow\Gr_{n-2}^W(H).
\]

\paragraph{8.5.} If $X$ is a nonsingular projective variety over an algebraically closed field $k_0$, we define
\[
L:\Z_l(-1)\otimes H^*(X,\Z_l)\longrightarrow H^*(X,\Z_l)
\]
as cup product with the cohomology class of a hyperplane section. One notes a formal analogy between $L$ and $M$; just as $M$ is defined by an action of $\Z_l(1)$, one may regard $L$ as defined by an action of $\Z_l(-1)$; $L$ raises the degree by $2$, and $\Gr M$ (8.4) lowers it by $2$.

\paragraph{9.} Let $D$ be the unit disk, $D^*=D-\{0\}$ and $X$
\[
\begin{array}{ccc}
X&\lhook\joinrel\longrightarrow&\mathbb P^r(\C)\times D\\[-0.2em]
{}^{f}\searrow&&\swarrow^{\pr_2}\\[-0.2em]
& D&
\end{array}
\]
a family of projective varieties parametrized by $D$, with $f$ proper and $f|D^*$ smooth. Keep the notation of 8, and recall that in the analogy between a henselian trait and a small neighborhood of $0$ in the complex line, we have the following dictionary (note that the spectrum of the ring of germs of holomorphic functions at $0$ is a henselian trait):
\[
\tag*{9.1.}
\begin{array}{r@{\quad\vrule\quad}l}
D & \Spec(V)\\
D^* & \Spec(K)\\
\text{a universal covering }\wt D^*\text{ of }D^* & \Spec(\barK)\\
\text{fundamental group }\pi_1(D^*) & \text{inertia group }I\\
\text{(with }\pi_1(D^*)=\Z\simeq\Z(1)_\Z\text{)} & \text{(with }I_l=\Z_l(1)\text{)}\\
X & \text{projective scheme }X\text{ over }\Spec(V)\\
X^*=f^{-1}(D^*) & X_K\\
\wt X=X\times_D\wt D^* & X_{\barK}\\
\text{local system }R^if_*\Z|D^* & \text{Galois module }H^i(X_{\barK},\Z_l)\\
H^i(\wt X,\Z) & H^i(X_K,\Z_l) % [sic] printed H^i(X_K,\Z_l); intended H^i(X_{\barK},\Z_l).
\end{array}
\]
Note that $\wt X$ is homotopy equivalent to each of the fibers $X_t=f^{-1}(t)$ ($t\in D^*$): $H^i(X_{\barK},\Z_l)$ again has $H^i(X_t,\Z)$ as its analogue, and the action of $I$ corresponds to the monodromy transformation $T$.

Here again, it is known that a subgroup of finite index in $\pi_1(D^*)$ acts unipotently on $H^i(\wt X,\Q)=H^i(X_t,\Q)$. Let us place ourselves in the semistable case where all of $\pi_1(D^*)$ acts unipotently (this amounts to replacing $D$ by a finite covering), and let $T$ be the action of the canonical generator of $\pi_1(D^*)$.

By 3 and 8, one expects that $H^i(\wt X,\Q)\simeq H^i(X_t,\Q)$ is equipped with an increasing filtration $W$, that $\Gr_n^W(H^i(\wt X,\Q)$ is equipped with a Hodge structure of weight $n$, that $\log T(W_n)\subset W_{n-2}$ and that $\log T$ induces a morphism of Hodge structures % [sic] printed \Gr_n^W(H^i(\wt X,\Q); intended \Gr_n^W(H^i(\wt X,\Q)).
\[
M_n:\Z(-1)\otimes\Gr_n^W(H^i)\longrightarrow\Gr_{n-2}^W(H^i). % [sic] printed \Z(-1); intended \Z(1).
\]
One would also like (8.2), and not only (8.3) and (8.4), to have an analogue.

In fact, one can define, for each vector $u$ in the tangent space to $D$ at $\{0\}$, a mixed Hodge structure $\calH_u$ on $H^i(\wt X,\Z)$. The filtration $W$ and the Hodge structures on the $\Gr_n^W(H^i)$ are independent of $u$, and the dependence of $\calH_u$ on $u$ is expressed simply in terms of $T$. In analogy with (8.2), one finds that, for every $u$, $\log T$ induces a homomorphism of mixed Hodge structures
\[
M:\Z(1)\otimes H^i(\wt X,\Z)\longrightarrow H^i(\wt X,\Z).
\]
Finally, the analogy 8.5 is not misleading (but here, the fact that $f|D^*$ is assumed to be proper and smooth is no doubt essential). One proves that
\[
(\log T)^k:\Gr_{n+k}^W(H^n(\wt X,\Q))\longrightarrow\Gr_{n-k}^W(H^n(\wt X,\Q))
\]
is an isomorphism for every $k$ (cf. [6], IV 6, cor. to th. 5). This characterizes the filtration $W$. So far, an analogue of Hodge's positivity theorem (cf. [6], IV 7, cor. to th. 7) is available only in very special cases. One hopes that the mixed structures $\calH_u$ determine the asymptotic behavior, as $t\to0$, of the family of pure structures $H^i(X_t,\Z)$ ($t\in D^*$).

\begin{center}
\bfseries BIBLIOGRAPHY
\end{center}
\begin{enumerate}[label={[\arabic*]},leftmargin=2.2em,itemsep=0.1em]
\item P. Deligne. — Hodge theory, to appear in \emph{Publ. Math. I. H. E. S.}, 40.
\item M. Demazure. — Motives of algebraic varieties, \emph{Sém. Bourbaki} (1969-1970), exp. 365.
\item J.-P. Serre. — Kähler analogues of certain conjectures of Weil, \emph{Ann. of Math.}, 71, 2 (1960), pp. 392-394.
\item — and J. Tate. — Good reduction of abelian varieties, \emph{Ann. of Math.}, 88, 3 (1968), pp. 492-517.
\item A. Weil. — Number of solutions of equations in finite fields, \emph{Bull. Amer. math. Soc.}, 55 (1949), pp. 497-508.
\item —. — Introduction to the study of Kähler manifolds, \emph{Act. Sci. et Ind.}, 1267, Hermann (1958).
\end{enumerate}

\begin{flushleft}
I. H. E. S.\\
35, route de Chartres,\\
91-Bures-sur-Yvette\\
(France)
\end{flushleft}

\end{document}
